
When practitioners choose a data attribution method to debug or understand their model, does the choice of method interact with the underlying architecture? Prior work has evaluated methods in isolation---EK-FAC on decoder-only LLMs \citep{grosse2023studying}, TRAK on BERT and CLIP separately \citep{park2023trak}---but no systematic comparison exists across architectures at matched conditions. This study addresses that gap by testing whether the answer depends on the method: specifically, whether TRAK is architecture-agnostic, TracIn favors encoders, and EK-FAC shows the theoretically expected decoder advantage.

Data attribution methods estimate the influence of training examples on model predictions. As foundation models grow in scale and deployment, understanding which training examples drive specific behaviors becomes important for debugging, data quality assessment, and model auditing. Three prominent methods---TRAK, EK-FAC, and TracIn---make fundamentally different approximations to the intractable influence function, trading computational cost for accuracy. However, prior evaluations have treated architecture as a fixed context rather than an experimental variable.

The hypothesis motivating this work is that attention structure should interact with attribution approximations. Encoder-only models like BERT use bidirectional attention, where each token attends to all others, creating dense gradient flow patterns. Decoder-only models like GPT-2 use causal attention masks that zero out the upper triangle of attention weights, concentrating gradient flow in the lower triangle. These structural differences propagate through backpropagation, affecting the Hessian curvature that approximation methods must capture or circumvent.

To test this hypothesis, matched experiments compare BERT-base and GPT-2 (both 12 layers, ~110-125M parameters) on SST-2 mislabeled detection with 5\% injected label noise. TRAK, EK-FAC, and TracIn are evaluated at three compute budgets each, measuring mislabeled detection AUC.

The findings reveal three patterns:

First, TRAK exhibits architecture-invariance, with less than 1\% AUC difference between BERT and GPT-2 at all compute budgets (0.36\%, 0.56\%, 0.11\% at projection dimensions 64, 256, 1024 respectively). Its random projection mechanism appears to average out architecture-specific gradient patterns.

Second, TracIn shows a directional BERT advantage (~3.03\% at 1 checkpoint, ~1.42\% at 2 checkpoints, ~1.13\% at 3 checkpoints). However, with only 2 random seeds, these differences are not statistically significant (p=0.26-0.39).

Third, EK-FAC does not show the predicted GPT-2 advantage. Despite theoretical expectations that Kronecker factorization fits causal structures better \citep{grosse2023studying}, only 0.27\% difference is observed (p=0.90), statistically indistinguishable from zero.

The causal mechanism underlying these patterns is validated through intermediate measurements. Attention sparsity differs by 98.82\% between architectures (BERT upper-triangle sparsity: 1.18\%, GPT-2: 100\%). This structural difference creates an $11\times$ gap in top Hessian eigenvalue (GPT-2: 0.502, BERT: 0.046), confirming that attention structure fundamentally shapes the loss landscape.

These results provide guidance for practitioners. For cross-architecture pipelines or when architecture may vary, TRAK is a robust default choice. For encoder-only work, TracIn may offer advantages. EK-FAC performs equivalently on both architectures, contrary to prior theoretical expectations.
